% Importado desde: 2026-03-28_Estructura_de_Simetrias_en_el_Par_Weyl.tex
\chapter{Derivacion Pedagogica: --- estructura de simetrias en el par Weyl controlado}
\label{ch:simetrias-par-weyl}






% verified-by: validators/chirality.py::test_canonical_weyl_det
% verified-by: validators/chirality.py::test_opposite_weyl_opposite_chirality
% verified-by: validators/chirality.py::test_parity_flips_chirality
% verified-by: validators/chirality.py::test_three_sigma_trace_levi_civita

\section{Objetivo}
Ya tenemos un par Weyl limpio, una lectura topologica de su quiralidad y una construccion discreta controlada. El siguiente paso natural es ordenar con mucho cuidado la estructura de simetrias.

La pregunta que guia este capitulo es:
\begin{displayquote}
\textbf{que simetrias tiene realmente el sistema en cada nivel del problema, y que parte de esa estructura podria ser relevante para una interpretacion de tipo espacio-tiempo emergente?}
\end{displayquote}

Este punto es importante porque en nuestro camino aparecen varios niveles distintos:
\begin{enumerate}[label=\arabic*.]
    \item simetrias microscopicas de red;
    \item simetrias del Hamiltoniano de Bloch;
    \item simetrias efectivas del sector Dirac ligero;
    \item simetrias remanentes despues del desdoblamiento a Weyl;
    \item simetrias geometricas y topologicas asociadas a la curvatura de Berry.
\end{enumerate}

Si no distinguimos esos niveles, corremos el riesgo de confundir:
\begin{enumerate}[label=\arabic*.]
    \item una simetria exacta de la microfisica;
    \item una simetria solo aproximada del infrarrojo;
    \item una propiedad topologica que no es, estrictamente, una simetria.
\end{enumerate}

\section{El modelo concreto que usamos}
Seguimos con el Hamiltoniano ya construido:
\begin{equation}
H(\bm{k})=
\frac{v}{a}\sum_{i=1}^3 \sin(k_i a)\,\alpha_i
+
M(\bm{k})\,\beta
+
b\,\mathbb{I}_\tau\otimes \sigma_z,
\label{c11:eq:H}
\end{equation}
con
\begin{equation}
M(\bm{k})=
m+r\sum_{i=1}^3\left(1-\cos(k_i a)\right),
\label{c11:eq:M}
\end{equation}
y
\begin{equation}
\alpha_i=\tau_z\otimes \sigma_i,
\qquad
\beta=\tau_x\otimes \mathbb{I}_2.
\label{c11:eq:ab}
\end{equation}

El termino proporcional a \(b\) es el que separa el sector Dirac ligero en un par Weyl a lo largo del eje \(z\).

\section{Primera capa: simetrias microscopicas de red}
\subsection{Traslaciones discretas}
La simetria mas basica de una red periodica es la invariancia bajo traslaciones discretas:
\begin{equation}
\bm{x}\mapsto \bm{x}+\bm{R},
\end{equation}
donde \(\bm{R}\) es un vector de la red.

En espacio recíproco, esta simetria es la razon por la cual el Hamiltoniano se organiza como \(H(\bm{k})\) y \(\bm{k}\) vive en una zona de Brillouin compacta.

Esta simetria:
\begin{enumerate}[label=\arabic*.]
    \item es exacta en el modelo de red;
    \item no es una traslacion continua del espacio fisico;
    \item es la base tanto del analisis de bandas como del problema de doblamiento.
\end{enumerate}

\subsection{Grupo puntual}
Sin el termino \(b\), el modelo de Wilson--Dirac de referencia tiene una estructura espacial isotropa en \(x,y,z\), al menos a nivel del Hamiltoniano escrito con los tres ejes tratados simetricamente.

Pero al encender
\begin{equation}
b\,\mathbb{I}\otimes \sigma_z,
\end{equation}
el eje \(z\) queda privilegiado. Eso reduce el grupo puntual efectivo:
\begin{enumerate}[label=\arabic*.]
    \item ya no hay simetria completa que intercambie \(x,y,z\);
    \item queda distinguido el subgrupo de transformaciones que preservan el eje \(z\);
    \item en el limite continuo isotropo, eso se ve como una simetria de rotacion alrededor de \(z\).
\end{enumerate}

\begin{figure}[h!]
\centering
\begin{tikzpicture}[
    box/.style={draw, rounded corners, thick, align=center, minimum width=4cm, minimum height=1.3cm},
    arr/.style={-{Latex[length=3mm]}, thick}
]
\node[box, fill=blue!7] (a) at (0,0) {red discreta\\traslaciones + grupo puntual};
\node[box, fill=yellow!10] (b) at (5.2,0) {perturbacion\\\(b\,\sigma_z\)};
\node[box, fill=green!7] (c) at (10.4,0) {eje \(z\) privilegiado\\simetria reducida};
\draw[arr] (a) -- (b);
\draw[arr] (b) -- (c);
\end{tikzpicture}
\caption{El termino que crea el par Weyl tambien reduce la simetria espacial del modelo.}
\end{figure}

\section{Segunda capa: inversion y reversal temporal}
Estas dos simetrias son centrales porque, en materiales tipo Weyl, romper una de ellas suele ser el mecanismo que desdobla un punto Dirac.

\subsection{Inversion}
Definimos el operador
\begin{equation}
\mathcal{P}=\beta=\tau_x\otimes \mathbb{I}_2.
\label{c11:eq:P}
\end{equation}

Queremos verificar si
\begin{equation}
\mathcal{P}H(\bm{k})\mathcal{P}^{-1}=H(-\bm{k}).
\label{c11:eq:Pinv}
\end{equation}

Usando \(\beta\alpha_i\beta=-\alpha_i\), \(\beta^2=\mathbb{I}\) y el hecho de que \(\beta\) conmuta con \(\mathbb{I}\otimes \sigma_z\), obtenemos:
\begin{align}
\mathcal{P}H(\bm{k})\mathcal{P}^{-1}
&=
\frac{v}{a}\sum_i \sin(k_i a)(-\alpha_i)
+
M(\bm{k})\beta
+
b\,\mathbb{I}\otimes \sigma_z
\\
&=
\frac{v}{a}\sum_i \sin(-k_i a)\alpha_i
+
M(-\bm{k})\beta
+
b\,\mathbb{I}\otimes \sigma_z
\\
&=
H(-\bm{k}).
\end{align}

Entonces la inversion se preserva.

\subsection{Reversal temporal}
Tomamos el operador de reversal temporal tipo espin \(1/2\):
\begin{equation}
\mathcal{T}=
\left(\mathbb{I}_\tau\otimes i\sigma_y\right)\mathcal{K},
\label{c11:eq:T}
\end{equation}
donde \(\mathcal{K}\) denota conjugacion compleja.

Este operador satisface
\begin{equation}
\mathcal{T}\sigma_i \mathcal{T}^{-1} = -\sigma_i.
\end{equation}
Por tanto:
\begin{enumerate}[label=\arabic*.]
    \item la parte cinetica del Hamiltoniano es invariante, porque \(\sin(k_i a)\) cambia de signo bajo \(\bm{k}\to -\bm{k}\) y \(\sigma_i\) tambien;
    \item el termino \(M(\bm{k})\beta\) es invariante;
    \item pero el termino \(b\,\mathbb{I}\otimes \sigma_z\) cambia de signo.
\end{enumerate}

Explicitamente,
\begin{equation}
\mathcal{T}\left(b\,\mathbb{I}\otimes \sigma_z\right)\mathcal{T}^{-1}
=
-b\,\mathbb{I}\otimes \sigma_z.
\end{equation}

Asi que
\begin{equation}
\mathcal{T}H(\bm{k})\mathcal{T}^{-1}\neq H(-\bm{k})
\end{equation}
cuando \(b\neq 0\).

\subsection{Lectura fisica}
Tenemos entonces la situacion caracteristica:
\begin{equation}
\text{inversion preservada},
\qquad
\text{reversal temporal rota}.
\end{equation}

Esta es una de las rutas estandar para obtener un par Weyl. Y encaja exactamente con el comportamiento de nuestro modelo.

\section{Tercera capa: el sector Dirac ligero antes del desdoblamiento}
Si tomamos primero \(b=0\) y nos quedamos en el subsector central, el Hamiltoniano efectivo era
\begin{equation}
H_{\text{Dirac}}(\bm{q})=
v\sum_{i=1}^3 q_i\alpha_i + m\beta.
\label{c11:eq:Diraceff}
\end{equation}

Cuando \(v_x=v_y=v_z=v\), esta expresion tiene una forma isotropa. En ese caso:
\begin{enumerate}[label=\arabic*.]
    \item el espectro depende solo de \(\abs{\bm{q}}\);
    \item el sector ligero admite una lectura relativista efectiva;
    \item el algebra de Dirac y los generadores \(\Sigma^{\mu\nu}\) quedan bien definidos, como ya vimos.
\end{enumerate}

\subsection{Punto delicado}
Esta simetria relativista efectiva no era microscopica. Era una simetria emergente del infrarrojo:
\begin{enumerate}[label=\arabic*.]
    \item aproximada;
    \item dependiente del limite de baja energia;
    \item compatible con una microfisica de red que no es Lorentz-invariante.
\end{enumerate}

Este matiz es crucial para el proyecto de espacio-tiempo emergente.

\section{Cuarta capa: que cambia al pasar de Dirac a Weyl}
Al encender \(b\neq 0\), el sector ligero queda descrito por
\begin{equation}
H_{\text{light}}(\bm{q})=
v\sum_{i=1}^3 q_i \alpha_i
+
m\beta
+
b\,\mathbb{I}\otimes \sigma_z.
\label{c11:eq:light}
\end{equation}

Este Hamiltoniano ya no tiene la misma simetria que \eqref{c11:eq:Diraceff}. El vector axial \(b\) introduce una direccion espacial fija.

\subsection{Lectura como campo de fondo}
Una forma muy util de entender esto es tratar a \(b\) como un fondo axial fijo. Entonces:
\begin{enumerate}[label=\arabic*.]
    \item las ecuaciones pueden seguir escribiendose con objetos tipo Dirac;
    \item pero el fondo selecciona una direccion privilegiada;
    \item por tanto la simetria efectiva completa ya no es \(SO(3,1)\), sino un subgrupo que preserva esa estructura.
\end{enumerate}

\subsection{Consecuencia}
Esto nos obliga a distinguir entre dos nociones:
\begin{enumerate}[label=\arabic*.]
    \item \textbf{covariancia formal:} podemos escribir el modelo con matrices gamma y espinores;
    \item \textbf{simetria efectiva real:} el sistema fisico concreto ya no tiene toda la isotropia ni toda la Lorentz efectiva del caso Dirac no perturbado.
\end{enumerate}

\section{Quinta capa: simetria remanente del par Weyl}
Despues del desdoblamiento, los nodos quedan en
\begin{equation}
\bm{q}_\pm=(0,0,\pm q_0),
\end{equation}
con quiralidades opuestas.

\subsection{Inversion y el intercambio de nodos}
Como la inversion envia \(\bm{q}\to -\bm{q}\), intercambia los dos nodos:
\begin{equation}
\mathcal{P}:\quad \bm{q}_+ \leftrightarrow \bm{q}_-.
\end{equation}

Esto es importante porque muestra que los nodos no son independientes por completo: forman un par relacionado por simetria.

\subsection{Rotaciones alrededor de \(z\)}
Si localmente el Hamiltoniano alrededor de cada nodo es
\begin{equation}
H_\pm(\delta \bm{q})
\simeq
v\,\sigma_x \delta q_x
+
v\,\sigma_y \delta q_y
\pm
v_z\,\sigma_z \delta q_z,
\label{c11:eq:weyllocal}
\end{equation}
entonces:
\begin{enumerate}[label=\arabic*.]
    \item en el plano \(xy\) sigue habiendo simetria circular efectiva cuando \(v_x=v_y\);
    \item el eje \(z\) queda distinguido;
    \item la simetria continua efectiva local se reduce a rotaciones alrededor de \(z\), no a rotaciones completas en \(3D\).
\end{enumerate}

\subsection{Una simetria local distinta de la global}
Este punto merece subrayarse:
\begin{displayquote}
\textbf{cerca de cada nodo, el cono puede parecer casi relativista; pero el par completo lleva una separacion \(\pm q_0\hat{z}\) que rompe la simetria isotropa global del sector ligero.}
\end{displayquote}

Es decir:
\begin{enumerate}[label=\arabic*.]
    \item cada nodo por separado tiene una estructura muy simple;
    \item el conjunto de ambos nodos recuerda que existe un eje privilegiado;
    \item esa informacion global no puede ignorarse si uno piensa en una interpretacion tipo espacio-tiempo.
\end{enumerate}

\begin{figure}[h!]
\centering
\begin{tikzpicture}[
    box/.style={draw, rounded corners, thick, align=center, minimum width=4cm, minimum height=1.3cm},
    arr/.style={-{Latex[length=3mm]}, thick}
]
\node[box, fill=blue!7] (a) at (0,0) {Dirac ligero\\simetria efectiva mas alta};
\node[box, fill=yellow!10] (b) at (5.2,0) {fondo axial\\\(b\hat{z}\)};
\node[box, fill=green!7] (c) at (10.4,0) {par Weyl\\simetria reducida};
\draw[arr] (a) -- (b);
\draw[arr] (b) -- (c);
\end{tikzpicture}
\caption{El desdoblamiento no solo separa nodos: tambien reduce la simetria efectiva.}
\end{figure}

\section{Sexta capa: la curvatura de Berry no es una simetria}
En el capitulo anterior vimos que cada nodo Weyl carga flujo de curvatura de Berry y se comporta como un monopolo topologico.

Eso no debe confundirse con una simetria.

\subsection{Que es la curvatura de Berry}
La curvatura de Berry:
\begin{enumerate}[label=\arabic*.]
    \item es un objeto geometrico en el espacio de momentos;
    \item codifica la estructura de las autofunciones;
    \item lleva informacion topologica sobre las bandas.
\end{enumerate}

\subsection{Que no es}
No es:
\begin{enumerate}[label=\arabic*.]
    \item un grupo de transformaciones;
    \item un generador de invariancias dinamicas;
    \item una simetria en el mismo sentido que inversion o reversal temporal.
\end{enumerate}

\subsection{Como conviven ambas estructuras}
Lo correcto es decir:
\begin{enumerate}[label=\arabic*.]
    \item la estructura de simetrias determina que perturbaciones son posibles y que nodos pueden aparecer;
    \item una vez aparecidos los nodos, la curvatura de Berry describe su carga topologica;
    \item simetria y topologia estan relacionadas, pero no son lo mismo.
\end{enumerate}

\section{Lo que esto sugiere para una conexion con espacio-tiempo}
Aqui es donde conviene ser muy cuidadosos.

\subsection{Lo prometedor}
Hay dos aspectos muy suggestivos para un programa de espacio-tiempo emergente:
\begin{enumerate}[label=\arabic*.]
    \item la aparicion de espinores, matrices gamma y una dinamica efectiva tipo relativista;
    \item la posibilidad de interpretar la direccion privilegiada y los fondos efectivos como parametros geometricos emergentes.
\end{enumerate}

\subsection{Lo que obliga a ser prudentes}
Pero tambien hay limites serios:
\begin{enumerate}[label=\arabic*.]
    \item el par Weyl no tiene ya simetria Lorentz efectiva completa;
    \item existe un eje privilegiado \(\hat{z}\);
    \item la microfisica sigue siendo de red y, por tanto, anisotropa y discreta;
    \item la geometria de Berry vive en espacio de momentos, no directamente en espacio-tiempo fisico.
\end{enumerate}

\subsection{Consecuencia metodologica}
Si el objetivo final es una lectura \enquote{intrinseca} de espacio-tiempo, entonces el caso Weyl no debe tomarse como el modelo final, sino como una pieza intermedia muy instructiva:
\begin{enumerate}[label=\arabic*.]
    \item muestra como aparecen espinores y quiralidad;
    \item muestra como entran fondos que rompen simetrias;
    \item pero tambien muestra que no toda fase interesante topologicamente es automaticamente un buen candidato de espacio-tiempo emergente isotropo.
\end{enumerate}

\section{Una palabra sobre arcos de Fermi y estados de borde}
Tu intuicion de ponerlos en segundo plano aqui es razonable.

\subsection{Por que no son el centro del argumento actual}
Los arcos de Fermi y los estados de borde:
\begin{enumerate}[label=\arabic*.]
    \item dependen de la presencia de una frontera;
    \item son observables de borde;
    \item revelan la no trivialidad topologica del bulk.
\end{enumerate}

Pero si lo que buscamos es una conexion con un \enquote{espacio-tiempo intrinseco}, lo primero es entender el bulk y su algebra de simetrias, no la fisica de superficies.

\subsection{Como podrian volver mas adelante}
Eso no significa que no importen nunca. Podrian ser utiles despues como:
\begin{enumerate}[label=\arabic*.]
    \item diagnostico de la estructura topologica del modelo;
    \item pista sobre como aparecen grados de libertad efectivos en fronteras;
    \item analogia para horizontes, defectos o interfaces en modelos mas geometricos.
\end{enumerate}

Pero no son la pieza central del argumento de espacio-tiempo en esta etapa.

\section{Sintesis por niveles}
Ya podemos resumir la estructura de simetrias en una tabla conceptual:

\begin{center}
\begin{tabular}{p{3.5cm} p{4.4cm} p{5.3cm}}
\toprule
Nivel & Simetria principal & Comentario \\
\midrule
Red microscopica & traslaciones discretas + grupo puntual & exactas, no relativistas \\
Wilson--Dirac sin \(b\) & inversion, reversal temporal, sector Dirac ligero & simetria efectiva mas alta en el infrarrojo \\
Con \(b\neq 0\) & inversion preservada, reversal temporal rota & aparece eje privilegiado \\
Par Weyl local & rotaciones efectivas alrededor de \(z\), intercambio por inversion & dos nodos de quiralidad opuesta \\
Geometria de Berry & carga topologica monopolar & no es simetria, sino estructura geometrica \\
\bottomrule
\end{tabular}
\end{center}

\section{Mapa del siguiente paso}
Ahora el hilo natural puede seguir por dos rutas, y ambas ya estan mejor fundamentadas gracias a esta aclaracion:
\begin{enumerate}[label=\arabic*.]
    \item volver al sector Dirac isotropo y preguntar que condiciones hacen mas plausible una lectura tipo espacio-tiempo emergente;
    \item estudiar si el fondo axial \(b\hat{z}\) puede reinterpretarse como un campo geometrico efectivo, mas que como una simple perturbacion material.
\end{enumerate}

\section{Resumen final}
El par Weyl controlado que construimos no solo tiene una topologia interesante; tambien tiene una estructura de simetrias muy precisa. El modelo conserva inversion, rompe reversal temporal y selecciona un eje espacial privilegiado. Por eso, aunque cada nodo localmente se parezca mucho a un cono relativista, el sistema completo ya no posee la misma simetria efectiva que el sector Dirac isotropo original.

Este resultado es importante para el proyecto grande. Nos dice que la emergencia de espinores, quiralidad y carga topologica no implica automaticamente una buena candidata de espacio-tiempo isotropo. Tambien nos dice por que tiene sentido, por ahora, concentrarnos en la estructura bulk y en las simetrias antes que en arcos de Fermi o estados de borde: esos son diagnosticos valiosos, pero la pregunta de fondo sobre espacio-tiempo emergente se juega primero en la algebra y en la geometria efectiva del volumen.
